\documentclass[10pt,a4paper,twoside]{article}

\usepackage[utf8]{inputenc}
\usepackage[brazil,portuges]{babel}
\usepackage{csquotes}

\usepackage[
	backend=biber,
	style=abnt,
	language=auto,
]{biblatex}

\addbibresource{../refs/referencias.bib}
\usepackage{soul}
\usepackage{xcolor}
\usepackage{framed}

\usepackage{kpfonts}
\usepackage[T1]{fontenc}

\usepackage{a4wide}
    \setlength{\headheight}{12.5pt}

\usepackage{tabularx}

	\newcolumntype{C}{>{\centering\arraybackslash}X}

\usepackage{graphicx}
\usepackage{multirow}
\usepackage{multicol}
\usepackage{hhline}
\usepackage{booktabs}
\usepackage[inline]{enumitem}
\usepackage{mathtools}

	\newcommand{\Binary}{\ensuremath{\mathbb{B}}}
	\newcommand{\Natural}{\ensuremath{\mathbb{N}}}
	\newcommand{\Integer}{\ensuremath{\mathbb{Z}}}
	\newcommand{\Rational}{\ensuremath{\mathbb{Q}}}
	\newcommand{\Real}{\ensuremath{\mathbb{R}}}
	
	\DeclarePairedDelimiter{\floor}{\lfloor}{\rfloor}
	\DeclarePairedDelimiter{\ceil}{\lceil}{\rceil}
	\DeclarePairedDelimiter{\abs}{\lvert}{\rvert}
	
\usepackage{amsthm}
	
	\newtheorem{proposition}{Proposição}[section]
	\theoremstyle{definition}
	\newtheorem{definition}{Definição}[section]

\usepackage{multicol}
	\newcommand{\finishcols}{\vfill\vphantom{.}}

\usepackage[portuguese,algoruled,vlined,noend]{algorithm2e}

	\SetKwInput{Algoritmo}{Algoritmo}
	\SetKw{Retorne}{retorne}
	\SetKwFor{Para}{para}{faça}{fimpara}
	\SetKw{Escreva}{escreva}
	\DontPrintSemicolon

\usepackage{hyperref}
\usepackage[brazil]{varioref}
\usepackage[brazilian]{cleveref}

	\def\disciplinacodigotext{QXD0115}
	\def\disciplinanometext{Estrutura de Dados Avançada}
	\def\semestretext{2022.1}
	\def\provatype{Notas de Aula}
	\def\provatitle{Árvores AVL são balanceadas}
	\def\provatext{\provatype{} \ifx\provatype\empty {} \else {--} \fi \provatitle{}}
	\def\provacodtext{}
	\def\professorabbrv{Prof.}
	\def\professornometext{Arthur Araruna}

\usepackage{fancyhdr}

	\pagestyle{fancy}
	\fancyhead{}
	\lhead{\disciplinacodigotext}
	\chead{\provatext}
	\rhead{\semestretext}
	\fancyfoot{}
	\fancyfoot[LE,RO]{\colorbox{black!20}{\thepage}}

\begin{document}
\noindent
\begin{minipage}{\textwidth}
	\begin{tabularx}{\textwidth}{cXc}
		\multirow{4}{*}{\includegraphics[height=1.8cm]{brasaoufc.pdf}} & {\bf UNIVERSIDADE FEDERAL DO CEARÁ}         & \multirow{3}{*}{} \\
		                                                               & {\bf Campus de Quixadá}                     &                   \\
		                                                               & \professorabbrv\enspace\professornometext   &                   \\
		                                                               & \disciplinacodigotext - \disciplinanometext & {}                \\
	\end{tabularx}
\end{minipage}

\bigskip
\bigskip

\part*{\provatitle}

\section{Conceitos}

Vamos relembrar o nossos conceitos que serão necessários para o resultado a ser demonstrado:

\begin{definition}
	Uma árvore binária \( T \), com \( n \) nós, tal que \( h(T) = O(\lg n) \) é dita \emph{balanceada}.
\end{definition}

\begin{definition}
    O \emph{balanço} de um nó binário \( n \) é definido pela função \( bal(n) \coloneqq h(n.dir) - h(n.esq) \).
\end{definition}

\begin{definition}
    Uma árvore binária \( T \) é \emph{AVL} quando, para todo nó \( n \in T \), \( \abs{bal(n)} \leq 1 \).
\end{definition}

\section{Demonstração}

% TODO: Reescrever, pois não é exatamente a forma que encontramos no final.
Note que devemos mostrar que \( h \leq c \cdot \lg n \) para algum \( c > 0 \), com \( h \) sendo a altura de uma árvore AVL \ul{qualquer} e \( n \) sendo seu tamanho. Entretanto, nossa estratégia será um pouco mais restrita, mas permitirá concluirmos a mesma coisa.

Em vez de analisarmos uma árvore qualquer, vamos observar, para um dado \( h \) arbitrário, uma árvore AVL com altura \( h \) que contenha o mínimo possível de nós. Se essa árvore respeitar a desigualdade anterior, poderemos concluir que \emph{todas} as AVL com essa mesma altura também a respeitarão (já que para essa tal árvore o lado direito da desigualdade é o menor possível, todas as outras apresentarão valores ainda maiores, o que mantém a desigualdade).

\subsection{Estrutura da árvore}
% TODO: Reescrever. Ficou um pouco confuso por ter uma sentença muito grande e misturar "árvore anterior" e "sub-árvore anterior"...
Para entendermos como construir uma árvore como desejamos, veja que sempre que, em uma árvore AVL \( T \) com altura \( h \), encontrarmos um nó \( n \) com \( bal(n) = 0 \), poderíamos substituir uma de suas sub-árvores por outra AVL com uma altura 1 unidade menor que a que foi retirada sem invalidar a propriedade AVL de \( T \). Isso constrói uma nova árvore \( T^\prime \) com a mesma altura que a anterior \( T \) e certamente tamanho menor. Dessa forma, podemos concluir que a árvore desejada não pode ter nenhum nó com balanço nulo.

\subsection{Contagem de nós da árvore}
Vamos definir a função \( N(h) \) de forma a representar o menor número de nós possível para uma árvore AVL de altura \( h \), que chamaremos de \( T \).

Como vimos anteriormente, as sub-árvores de \( T \) devem possuir alturas \( h-1 \) e \( h-2 \), não necessariamente nesta ordem. Além disso, podemos afirmar que ambas se tratam de árvores AVL com o mínimo de nós para a respectiva altura, caso contrário não faria sentido \( T \) possuir tamanho mínimo (já que elas poderiam ser trocadas por outras árvores de mesma altura e menos nós).

Dessa forma, podemos afirmar que o valor de \( N(h) \) pode ser calculado como \[ N(h) = \begin{cases}
    N(h-1) + N(h-2) + 1, & \text{se } h \geq 2; \\
    1                  , & \text{se } h = 1; \\
    0                  , & \text{se } h = 0.
\end{cases} \]

Veja que existe uma enorme semelhança entre a função \( N \) e a função \( F \) a seguir, que define a \emph{Sequência de Fibonacci}: \[ F(h) = \begin{cases}
    F(h-1) + F(h-2) , & \text{se } h \geq 2; \\
    1               , & \text{se } h = 1; \\
    0               , & \text{se } h = 0.
\end{cases} \]

De fato, para qualquer valor de \( h \), é o caso que \( N(h) \geq F(h) \). Vamos usar o valor \( f_h \) do \( h \)-ésimo termo dessa sequência para obtermos uma relação entre \( h \) e o número de nós \( N(h) \) da árvore, que denotaremos simplesmente por \( n \). Podemos calcular esse termo de forma não-recursiva como\footnote{Veja \url{https://pt.wikipedia.org/wiki/Sequência_de_Fibonacci\#Fórmula_explícita}} \[ f_h = \frac{1}{\sqrt{5}}\left( \frac{1 + \sqrt{5}}{2} \right)^h - \frac{1}{\sqrt{5}}\left( \frac{1 - \sqrt{5}}{2} \right)^h. \]

Como \( h > 0 \), podemos concluir\footnote{Ver apêndice para a argumentação mais detalhada.} que \[ \frac{1}{\sqrt{5}}\left( \frac{1 - \sqrt{5}}{2} \right)^h \leq 1 \implies f_h \geq \frac{1}{\sqrt{5}}\left( \frac{1 + \sqrt{5}}{2} \right)^h - 1. \]

\section{Conclusão}
Daí, considerando \( a = \frac{1 + \sqrt{5}}{2} \) e lembrando que \( n = N(h) \geq F(h) = f_h \), podemos, ao reescrever a desigualdade anterior, concluir que
\begin{eqnarray*} 
	n \geq f_h \geq & \displaystyle\frac{1}{\sqrt{5}} a^h - 1 & \implies \\
	n + 1 \geq & \displaystyle\frac{1}{\sqrt{5}} a^h & \implies \\
	\log_a (n + 1) \geq & h - \log_a \sqrt{5} & \implies \\
	\log_a (n + 1) + \log_a \sqrt{5} \geq & h.
\end{eqnarray*}

Isso significa que \( h \leq \log_a (n + 1) + \log_a \sqrt{5} = O(\lg n) \), como desejávamos.

\part*{Apêndice}

\end{document}